% !TeX program = XeLaTeX
% !TeX root = ../../shloka.tex
% --meta--
% deity: Svadha
% composer: Brahma
% chandas: Anushtup
% source: Devi Bhagavatam
% vakta: Narayana
% shrota: Narada
% wiki_title: Svadha Stotram
% --end-meta--

\sect{स्वधास्तोत्रम्}
\addtocounter{shlokacount}{26}

\uvacha{श्रीनारायण उवाच}

\twolineshloka
{स्वधोच्चारणमात्रेण तीर्थस्नायी भवेन्नरः}
{मुच्यते सर्वपापेभ्यो वाजपेयफलं लभेत्} % २७

\twolineshloka
{स्वधा स्वधा स्वधेत्येवं यदि वारत्रयं स्मरेत्}
{श्राद्धस्य फलमाप्नोति बलेश्च तर्पणस्य च} % २८

\twolineshloka
{श्राद्धकाले स्वधास्तोत्रं यः शृणोति समाहितः}
{स लभेच्छ्राद्धसम्भूतं फलमेव न संशयः} % २९

\twolineshloka
{स्वधा स्वधा स्वधेत्येवं त्रिसन्ध्यं यः पठेन्नरः}
{प्रियां विनीतां स लभेत्साध्वीं पुत्रगुणान्विताम्} % ३०

\twolineshloka
{पितॄणां प्राणतुल्या त्वं द्विजजीवनरूपिणी}
{श्राद्धाधिष्ठातृदेवी च श्राद्धादीनां फलप्रदा} % ३१

\twolineshloka
{नित्या त्वं सत्यरूपासि पुण्यरूपासि सुव्रते}
{आविर्भावतिरोभावौ सृष्टौ च प्रलये तव} % ३२

\twolineshloka
{ॐ स्वस्तिश्च नमः स्वाहा स्वधा त्वं दक्षिणा तथा}
{निरूपिताश्चतुर्वेदैः प्रशस्ताः कर्मिणां पुनः} % ३३

\twolineshloka
{कर्मपूर्त्यर्थमेवैता ईश्वरेण विनिर्मिताः}
{इत्येवमुक्त्वा स ब्रह्मा ब्रह्मलोके स्वसंसदि} % ३४

\twolineshloka
{तस्थौ च सहसा सद्यः स्वधा साविर्बभूव ह}
{तदा पितृभ्यः प्रददौ तामेव कमलाननाम्} % ३५

\threelineshloka
{तां सम्प्राप्य ययुस्ते च पितरश्च प्रहर्षिताः}
{स्वधास्तोत्रमिदं पुण्यं यः शृणोति समाहितः}
{स स्नातः सर्वतीर्थेषु वाञ्छितं फलमाप्नुयात्} % ३६

॥इति श्रीमद्देवीभागवते महापुराणेऽष्टादशसाहस्र्यां संहितायां नवमस्कन्धे नारायणनारदसंवादे स्वधोपाख्यानवर्णनं नाम चतुश्चत्वारिंशोऽध्यायः॥

\closesub
